% (User input window)
% 2026-02-16 23:16:42（Asia/Singapore）

\paragraph{最大熵马尔可夫链、对称化谱与小失真协同}
沿用定理 \ref{thm:pom-diagonal-rate-explicit-param-and-derivative} 与定理 \ref{thm:pom-diagonal-rate-small-distortion-hellinger-endpoint} 的记号。
本段把对角质量约束的最优耦合 \(P^\star_\delta\) 等价重写为“给定平稳分布与全局自环率的最大熵马尔可夫链”，
并进一步给出其可逆化对称谱的“对角 \(+\) 秩一”结构、世俗方程与端点协同的张量同态律。

\begin{theorem}[最大熵马尔可夫链：给定平稳分布与全局自环率时的唯一最优核]\label{thm:pom-maxent-markov-unique-optimal-kernel}
令 \(X\) 为有限集，取全支撑分布 \(w\in\Delta(X)\)。
对任意马尔可夫核 \(K(\cdot\mid x)\)（\(x\in X\)）考虑约束
\[
\sum_{x\in X}w(x)K(y\mid x)=w(y)\quad(\forall y\in X),
\qquad
\sum_{x\in X}w(x)K(x\mid x)\ge 1-\delta,\qquad \delta\in[0,1].
\]
定义熵率
\[
h(K):=\sum_{x\in X}w(x)\,H\!\bigl(K(\cdot\mid x)\bigr),
\qquad
H\!\bigl(K(\cdot\mid x)\bigr):=-\sum_{y\in X}K(y\mid x)\log K(y\mid x).
\]
则在上述可行集合上，\(h(K)\) 存在唯一最大化解 \(K^\star_\delta\)。
令耦合
\[
P^\star_\delta(x,y):=w(x)K^\star_\delta(y\mid x),
\]
则 \(P^\star_\delta\) 恰为
\[
R_w(\delta)=\inf\{I_P(X;Y):\ P_X=P_Y=w,\ \mathbb P_P(X=Y)\ge 1-\delta\}
\]
的唯一极小耦合，并满足恒等式
\[
\boxed{\;
R_w(\delta)=H(w)-h\!\bigl(K^\star_\delta\bigr)\; }.
\]
\end{theorem}
\begin{proof}
在平稳性约束下，\(P(x,y):=w(x)K(y\mid x)\) 满足 \(P_X=w\) 且
\(
P_Y(y)=\sum_x w(x)K(y\mid x)=w(y)
\)，故 \(P_X=P_Y=w\)。
同时
\(
\mathbb P_P(X=Y)=\sum_x w(x)K(x\mid x)
\)。
并且
\(
h(K)=H_P(Y\mid X)
\)。
因此最大化 \(h(K)\) 等价于在 \(\{P_X=P_Y=w,\ \mathbb P(X=Y)\ge 1-\delta\}\) 上最大化 \(H_P(Y\mid X)\)。
又由恒等式
\(
I_P(X;Y)=H(w)-H_P(Y\mid X)
\)，该最大化等价于最小化互信息 \(I_P(X;Y)\)，其最优值即为 \(R_w(\delta)\)。
由于 \(K\mapsto h(K)\) 在每个行向量上严格凹、且约束为线性，最大化器唯一；
等价地，\(P\mapsto I_P(X;Y)\) 在可行凸集上严格凸，极小化器亦唯一。
两者互为同一对象，从而得到结论。证毕。
\end{proof}

\begin{theorem}[对称化谱的“对角 \(+\) 秩一”结构与世俗方程]\label{thm:pom-maxent-markov-diagonal-plus-rankone-spectrum}
令 \(P^\star_\delta\) 与 \(K^\star_\delta\) 如定理 \ref{thm:pom-maxent-markov-unique-optimal-kernel}，
并记 \(D:=\mathrm{diag}(w)\)。
则存在唯一参数 \(\kappa_\delta\ge 0\) 与向量 \(u_\delta\in\RR_{>0}^{X}\)（唯一到整体符号，且可取 \(u_\delta\succ0\)）
使得
\[
\boxed{\;
P^\star_\delta=u_\delta u_\delta^{\mathsf T}+\kappa_\delta\,\mathrm{diag}(u_\delta^2)\;}
\]
亦即对 \(x\ne y\) 有 \(P^\star_\delta(x,y)=u_\delta(x)u_\delta(y)\)，并且
\(
P^\star_\delta(x,x)=(1+\kappa_\delta)u_\delta(x)^2
\)。
定义对称化算子
\[
S_\delta:=D^{-1/2}P^\star_\delta D^{-1/2},
\qquad
a_\delta:=D^{-1/2}u_\delta,
\qquad
d_x:=\kappa_\delta\,a_\delta(x)^2,
\]
则有严格分解
\[
\boxed{\;
S_\delta=a_\delta a_\delta^{\mathsf T}+\kappa_\delta\,\mathrm{diag}(a_\delta^2)
=a_\delta a_\delta^{\mathsf T}+\mathrm{diag}(d_x)\;}
\]
从而 \(S_\delta\succeq 0\) 且其全部特征值为非负实数。
并且 \(S_\delta\) 与 \(K^\star_\delta\) 在 \(\ell^2(w)\) 上相似，故二者具有相同实谱；
其中最大特征值为 \(\lambda_1=1\)。
\par
设 \(\lambda\in\CC\)。则 \(S_\delta\) 的特征多项式满足闭式
\[
\boxed{\;
\det(S_\delta-\lambda I)
=\Bigl(\prod_{x\in X}(d_x-\lambda)\Bigr)\Bigl(1+\sum_{x\in X}\frac{a_\delta(x)^2}{d_x-\lambda}\Bigr)\;}
\]
从而除去对角极点 \(d_x\) 外，全部特征值由一维世俗方程
\[
\boxed{\;
1+\sum_{x\in X}\frac{a_\delta(x)^2}{d_x-\lambda}=0\;}
\]
刻画，并满足标准的 interlacing 交错性质。
\end{theorem}
\begin{proof}
由定理 \ref{thm:pom-diagonal-rate-explicit-param-and-derivative} 的指数族形式，
对某组 \((\widehat u_\delta,\kappa_\delta,\widehat Z_\delta)\) 有
\(
P^\star_\delta(x,y)=\widehat u_\delta(x)\widehat u_\delta(y)\bigl(1+\kappa_\delta\mathbf 1_{\{x=y\}}\bigr)/\widehat Z_\delta
\)。
将 \(\widehat Z_\delta^{-1/2}\) 吸收进 \(\widehat u_\delta\) 得到 \(u_\delta:=\widehat u_\delta/\sqrt{\widehat Z_\delta}\)，即可写成
\(
P^\star_\delta=u_\delta u_\delta^{\mathsf T}+\kappa_\delta\mathrm{diag}(u_\delta^2)
\)；
由最优解唯一性可知 \(\kappa_\delta\) 与 \(u_\delta\)（取 \(u_\delta\succ0\)）唯一。
对称化分解由左右乘 \(D^{-1/2}\) 直接得到；两项均为正半定，故 \(S_\delta\succeq 0\)。
\par
由于 \(P^\star_\delta\mathbf 1=w\)，有
\[
S_\delta\sqrt w
=D^{-1/2}P^\star_\delta D^{-1/2}\sqrt w
=D^{-1/2}P^\star_\delta\mathbf 1
=D^{-1/2}w
=\sqrt w,
\]
从而 \(\lambda_1=1\)。
又 \(P^\star_\delta=D K^\star_\delta\)，故
\(
S_\delta=D^{1/2}K^\star_\delta D^{-1/2}
\)，给出相似与谱一致。
\par
最后，对 \(\mathrm{diag}(d_x)+a_\delta a_\delta^{\mathsf T}\) 应用矩阵行列式引理
\(
\det(A+uv^{\mathsf T})=\det(A)\bigl(1+v^{\mathsf T}A^{-1}u\bigr)
\)
即可得到特征多项式闭式与世俗方程。证毕。
\end{proof}

\begin{corollary}[最优耦合的二元潜变量表示]\label{cor:pom-maxent-markov-optimal-coupling-latent-binary-mixture}
沿用定理 \ref{thm:pom-maxent-markov-diagonal-plus-rankone-spectrum} 的 \(u_\delta,\kappa_\delta\)。
记
\[
A_\delta:=\sum_{x\in X}u_\delta(x),
\qquad
\theta_\delta:=\kappa_\delta\sum_{x\in X}u_\delta(x)^2,
\qquad
\pi_\delta(x):=\frac{u_\delta(x)}{A_\delta},
\qquad
\rho_\delta(x):=\frac{u_\delta(x)^2}{\sum_{y\in X}u_\delta(y)^2}.
\]
则由概率归一化恒等式 \(\sum_{x,y}P^\star_\delta(x,y)=1\) 推出
\[
A_\delta^2+\theta_\delta=1.
\]
并且最优耦合 \(P^\star_\delta\) 可写为两端混合：
\[
\boxed{\;
P^\star_\delta
=A_\delta^2\,(\pi_\delta\otimes\pi_\delta)
\;+\;\theta_\delta\,\mathrm{Diag}(\rho_\delta)\;}
\]
其中 \(\mathrm{Diag}(\rho_\delta)\) 表示集中在对角集合 \(\{(x,x)\}\) 的分布。
等价地，存在二元潜变量 \(J\in\{0,1\}\) 使得在 \(J=0\) 时 \((X,Y)\sim\pi_\delta\otimes\pi_\delta\) 独立同分布，在 \(J=1\) 时 \(X=Y\) 且 \(X\sim\rho_\delta\)，并且
\[
\mathbb P(J=0)=A_\delta^2,\qquad \mathbb P(J=1)=\theta_\delta,
\qquad
{\rm Law}(X,Y)=P^\star_\delta.
\]
\end{corollary}
\begin{proof}
由定理 \ref{thm:pom-maxent-markov-diagonal-plus-rankone-spectrum} 的分解 \(P^\star_\delta=u_\delta u_\delta^{\mathsf T}+\kappa_\delta\,\mathrm{diag}(u_\delta^2)\)，对总质量求和得
\[
1=\sum_{x,y}P^\star_\delta(x,y)=\Bigl(\sum_x u_\delta(x)\Bigr)^2+\kappa_\delta\sum_x u_\delta(x)^2=A_\delta^2+\theta_\delta.
\]
另一方面，
\(
u_\delta u_\delta^{\mathsf T}=A_\delta^2(\pi_\delta\otimes\pi_\delta)
\)
且
\(
\kappa_\delta\,\mathrm{diag}(u_\delta^2)=\theta_\delta\,\mathrm{Diag}(\rho_\delta)
\)，
从而得到混合表示。
将两部分视作对 \(J\) 的条件分布并按权重边缘化即可得到潜变量生成模型。证毕。
\end{proof}

\begin{corollary}[无振荡性与两端谱塌缩]\label{cor:pom-maxent-markov-no-oscillation-endpoints}
令 \(\delta_0:=1-p_2(w)=1-\sum_{x\in X}w(x)^2\)。
则对任意 \(\delta\in[0,\delta_0]\)，最优核 \(K^\star_\delta\) 在 \(\ell^2(w)\) 上的全部特征值均属于区间 \([0,1]\)。
特别地，\(K^\star_\delta\) 不存在负谱模态，因而在谱意义下不携带周期振荡成分。
\par
当 \(\delta=\delta_0\) 时，独立耦合 \(P=w\otimes w\) 可行且最优，因而
\[
\mathrm{Spec}(K^\star_{\delta_0})=\{1,0,\dots,0\}.
\]
当 \(\delta=0\) 时，对角耦合 \(P=\mathrm{Diag}(w)\) 最优，因而
\[
\mathrm{Spec}(K^\star_0)=\{1,1,\dots,1\}.
\]
\end{corollary}
\begin{proof}
由定理 \ref{thm:pom-maxent-markov-diagonal-plus-rankone-spectrum}，\(S_\delta\succeq0\) 且与 \(K^\star_\delta\) 相似，故 \(\mathrm{Spec}(K^\star_\delta)\subseteq[0,\infty)\)。
另一方面，\(K^\star_\delta\) 为马尔可夫核，算子范数满足 \(\|K^\star_\delta\|_{\ell^2(w)\to \ell^2(w)}\le 1\)，且 \(\lambda_1=1\)；
故全部特征值属于 \([0,1]\)。
\par
端点情形：当 \(\delta=\delta_0\) 时独立耦合满足 \(\mathbb P(X=Y)=p_2(w)=1-\delta_0\) 且互信息为零，故最优且
\(
K(y\mid x)=w(y)
\) 为秩一投影，谱为 \(\{1,0,\dots,0\}\)。
当 \(\delta=0\) 时约束强制 \(\mathbb P(X=Y)=1\)，唯一可行耦合为 \(\mathrm{Diag}(w)\)，相应 \(K\) 为恒等算子，谱为全 \(1\)。
证毕。
\end{proof}

\begin{corollary}[小失真谱门槛：第二特征值的首阶系数]\label{thm:pom-maxent-markov-small-distortion-lambda2-slope}
取全支撑分布 \(w\in\Delta(X)\) 且 \(|X|=n\ge 2\)。
令 \(K^\star_\delta\) 为定理 \ref{thm:pom-maxent-markov-unique-optimal-kernel} 的唯一最优核，并令
\(
0=\mu_1<\mu_2\le\cdots\le\mu_n
\)
为定理 \ref{thm:pom-diagonal-rate-critical-line-laplacian-first-order} 所定义之临界线拉普拉斯 \(L_w\) 的特征值。
则当 \(\delta\downarrow 0\) 时有
\[
\boxed{\;
1-\lambda_2(K^\star_\delta)=\mu_2\,\delta+o(\delta)\; }.
\]
并且 \(\mu_2\) 由定理 \ref{thm:pom-diagonal-rate-critical-spectrum-secular-interlacing} 的世俗方程与交错定位唯一确定。
\end{corollary}
\begin{proof}
由定理 \ref{thm:pom-diagonal-rate-critical-line-laplacian-first-order} 得到全谱首阶展开
\(
\lambda_i(K^\star_\delta)=1-\delta\mu_i+o(\delta)
\)
（其中 \(\lambda_1\equiv 1\)），取 \(i=2\) 即得。
世俗方程与定位为定理 \ref{thm:pom-diagonal-rate-critical-spectrum-secular-interlacing} 的直接推论。证毕。
\end{proof}

\begin{theorem}[小失真协同：乘积分布下的最优率严格优于任意坐标分离策略]\label{thm:pom-diagonal-rate-small-distortion-synergy-product}
取两有限集 \(X_1,X_2\)，分布 \(w\in\Delta(X_1)\)、\(v\in\Delta(X_2)\)，并令乘积分布 \(\omega:=w\otimes v\in\Delta(X_1\times X_2)\)。
记
\[
C_{1/2}(w):=A_{1/2}(w)^2-1,\qquad C_{1/2}(v):=A_{1/2}(v)^2-1,\qquad
C_{1/2}(\omega):=A_{1/2}(\omega)^2-1.
\]
则当 \(\delta\downarrow 0\) 时有小失真展开
\[
R_{\omega}(\delta)=H(w)+H(v)-\delta\log\frac{C_{1/2}(\omega)}{\delta}+O(\delta).
\]
\par
称“坐标分离策略”为取两耦合 \(P_1,P_2\) 分别作用于 \(X_1,X_2\)，并用乘积耦合 \(P=P_1\otimes P_2\) 生成 \(X_1\times X_2\) 的耦合。
若
\(
\mathbb P(X_1=Y_1)\ge 1-\delta_1
\)、
\(
\mathbb P(X_2=Y_2)\ge 1-\delta_2
\)，
则
\(
\mathbb P(Z=Z')\ge (1-\delta_1)(1-\delta_2)=1-(\delta_1\oplus\delta_2)
\)，其中 \(\delta_1\oplus\delta_2:=\delta_1+\delta_2-\delta_1\delta_2\)。
在该类策略中达到对角质量 \(1-\delta\) 的最优互信息满足
\[
\inf_{\delta_1\oplus\delta_2=\delta}\bigl(R_w(\delta_1)+R_v(\delta_2)\bigr)
=H(w)+H(v)-\delta\log\frac{C_{1/2}(w)+C_{1/2}(v)}{\delta}+O(\delta).
\]
并且存在严格协同增益：只要 \(C_{1/2}(w),C_{1/2}(v)>0\)（等价于两分布均非单点），就有
\[
\inf_{\delta_1\oplus\delta_2=\delta}\bigl(R_w(\delta_1)+R_v(\delta_2)\bigr)-R_{w\otimes v}(\delta)
=\delta\log\Bigl(1+\frac{C_{1/2}(w)C_{1/2}(v)}{C_{1/2}(w)+C_{1/2}(v)}\Bigr)+o(\delta)\;>\;0.
\]
\end{theorem}
\begin{proof}
由定理 \ref{thm:pom-diagonal-rate-small-distortion-hellinger-endpoint}，
\(
R_w(\delta)=H(w)-\delta\log\frac{C_{1/2}(w)}{\delta}+O(\delta)
\)，同理对 \(v\) 与 \(\omega\) 成立。
又 \(H(\omega)=H(w)+H(v)\) 且
\(
A_{1/2}(\omega)=A_{1/2}(w)A_{1/2}(v)
\)，故
\[
C_{1/2}(\omega)
=A_{1/2}(\omega)^2-1
=(1+C_{1/2}(w))(1+C_{1/2}(v))-1
=C_{1/2}(w)+C_{1/2}(v)+C_{1/2}(w)C_{1/2}(v).
\]
第一式由此立即得到。
\par
对分离类，约束 \(\delta_1\oplus\delta_2=\delta\) 等价于 \(\delta_1+\delta_2=\delta+O(\delta^2)\)，
故在 \(\delta\downarrow 0\) 的首阶上可将其替换为 \(\delta_1+\delta_2=\delta\)。
将端点展开代入并用拉格朗日乘子最大化
\(
\delta_1\log(C_{1/2}(w)/\delta_1)+\delta_2\log(C_{1/2}(v)/\delta_2)
\)
得到最优配比 \(\delta_1:\delta_2=C_{1/2}(w):C_{1/2}(v)\) 与系数 \(C_{1/2}(w)+C_{1/2}(v)\)，从而给出分离最优展开。
最后将 \(C_{1/2}(\omega)=C_{1/2}(w)+C_{1/2}(v)+C_{1/2}(w)C_{1/2}(v)\) 与上述系数比较即可得到严格差距公式。证毕。
\end{proof}

\begin{theorem}[\texorpdfstring{$k$}{k} 重协同放大：小失真端点的首阶系数按乘法增长]\label{thm:pom-diagonal-rate-small-distortion-synergy-k}
取 \(k\ge 2\) 个全支撑分布 \(w^{(1)},\dots,w^{(k)}\)，并令 \(\omega=\bigotimes_{i=1}^k w^{(i)}\)。
记 \(C_i:=C_{1/2}(w^{(i)})\) 以及
\[
C_{\mathrm{joint}}:=C_{1/2}(\omega)=\prod_{i=1}^k(1+C_i)-1,\qquad
C_{\mathrm{sep}}:=\sum_{i=1}^k C_i.
\]
则当 \(\delta\downarrow 0\) 时成立
\[
R_{\omega}(\delta)=\sum_{i=1}^k H\!\bigl(w^{(i)}\bigr)-\delta\log\frac{C_{\mathrm{joint}}}{\delta}+O(\delta),
\]
而任意坐标分离策略的最优首阶系数至多为 \(C_{\mathrm{sep}}\)，即
\[
\inf_{\delta_1\oplus\cdots\oplus\delta_k=\delta}\ \sum_{i=1}^k R_{w^{(i)}}(\delta_i)
=\sum_{i=1}^k H\!\bigl(w^{(i)}\bigr)-\delta\log\frac{C_{\mathrm{sep}}}{\delta}+O(\delta).
\]
因此首阶协同增益为
\[
\delta\log\frac{C_{\mathrm{joint}}}{C_{\mathrm{sep}}}+o(\delta),
\]
其在 \(C_i\) 为常数阶时随 \(k\) 可呈指数级增长。
\end{theorem}
\begin{proof}
由 \(A_{1/2}(\omega)=\prod_i A_{1/2}(w^{(i)})\) 得
\(
1+C_{1/2}(\omega)=\prod_i (1+C_{1/2}(w^{(i)}))
\)，
从而 \(C_{\mathrm{joint}}=\prod_i(1+C_i)-1\)。
其余与定理 \ref{thm:pom-diagonal-rate-small-distortion-synergy-product} 相同：对每个因子使用端点展开并用多元拉格朗日乘子得到分离最优配比 \(\delta_i\propto C_i\)，从而首阶系数为 \(C_{\mathrm{sep}}\)。
比较 \(C_{\mathrm{joint}}\) 与 \(C_{\mathrm{sep}}\) 即得增益。证毕。
\end{proof}

\begin{corollary}[端点协同生成元的张量同态]\label{cor:pom-diagonal-rate-small-distortion-synergy-generator-homomorphism}
对任意全支撑分布 \(w\) 定义
\[
\mathcal J(w):=\log\bigl(1+C_{1/2}(w)\bigr)=2\log A_{1/2}(w).
\]
则对任意两分布 \(w,v\) 成立严格等式
\[
\boxed{\;
\mathcal J(w\otimes v)=\mathcal J(w)+\mathcal J(v)\;}
\]
等价地，
\(
1+C_{1/2}(w\otimes v)=(1+C_{1/2}(w))(1+C_{1/2}(v))
\)。
并且在小失真极限 \(\delta\downarrow 0\) 下，
\[
H(w)-R_w(\delta)=\delta\log\frac{e^{\mathcal J(w)}-1}{\delta}+O(\delta),
\]
因此 \(\mathcal J\) 给出端点协同系数的可加生成元。
\end{corollary}
\begin{proof}
由 \(A_{1/2}(w\otimes v)=A_{1/2}(w)A_{1/2}(v)\) 得
\(
\mathcal J(w\otimes v)=2\log A_{1/2}(w\otimes v)=2\log A_{1/2}(w)+2\log A_{1/2}(v)
\)。
最后一式由定理 \ref{thm:pom-diagonal-rate-small-distortion-hellinger-endpoint} 的端点展开重写得到。证毕。
\end{proof}

